\section{Hardware}
\label{sec:hardware}

The build is two devices: the collar worn by the cat and the station plugged in at home. There is no
custom PCB. Each device is a Seeed Studio XIAO module in a 3D-printed shell, chosen for the small
footprint, integrated radios, hand-solderable castellated pads, and, on the collar, on-module
sensors. The full parts list is in \cref{sec:appendix-bom}; print files and assembly steps are in
\texttt{hardware/README.md}. This section summarises the design and its reasoning.

\subsection{Collar}

The collar is built around the Seeed XIAO nRF52840 Sense. The nRF52840 provides a Bluetooth Low
Energy radio and enough flash and RAM to run quantised neural-network inference on its Cortex-M4F.
Critically for Meowtion, the \emph{Sense} variant carries both sensors the system needs on the module
itself, a six-axis IMU (LSM6DS3TR-C) and a PDM MEMS microphone, so the collar needs no external
sensor breakout and stays small and light.

The collar runs on its own 1S LiPo cell (3.7\,V, 100\,mAh, with a built-in protection circuit)
soldered to the XIAO's BAT pads and charged through the module's onboard divider. The build is
completed by a 20\,mm hydrophobic PTFE filter membrane over the mic hole, PETG front and back shells,
a flexible TPU skin, and a 10\,mm quick-release safety strap. The full parts list is in
\cref{sec:appendix-bom}.

The IMU supplies the motion stream that is always classified; the microphone supplies the short audio
window used to disambiguate behaviours that look alike by motion but differ by sound
(\cref{sec:on-device-ai}). The microphone port sits behind the PTFE membrane in the front shell;
correct membrane placement matters for audio quality and water resistance and is documented in the
hardware guide.

\begin{figure}[ht]
  \centering
  \begin{subfigure}{0.31\textwidth}
    \centering
    \includegraphics[width=\linewidth]{exploded.png}
    \caption{Exploded view.}
    \label{fig:exploded}
  \end{subfigure}\hfill
  \begin{subfigure}{0.31\textwidth}
    \centering
    \includegraphics[width=\linewidth]{internals.jpg}
    \caption{Assembled internals.}
    \label{fig:internals}
  \end{subfigure}\hfill
  \begin{subfigure}{0.31\textwidth}
    \centering
    \includegraphics[width=\linewidth]{membrane-placement.jpg}
    \caption{PTFE membrane over the mic hole.}
    \label{fig:membrane}
  \end{subfigure}
  \caption{Collar and station hardware. Full parts, print files, and assembly steps are in
           \texttt{hardware/README.md}.}
  \label{fig:collar-hw}
\end{figure}

\subsection{Station}

The station is a Seeed XIAO ESP32-S3. The ESP32-S3 has both Wi-Fi and BLE, so a single module acts as
the BLE central that listens to the collar and the Wi-Fi client that reaches the cloud. Mains-powered
over USB-C, it scans and stays online continuously without the wearable's power constraints. It
carries no sensors of its own; its job is to hear the collar, capture training clips when asked, relay
results to the cloud, and push models to the collar. Its parts are listed in \cref{sec:appendix-bom}.

\subsection{Enclosure and assembly}

The front shell is shared between the collar and the station; the back shells and the collar's
flexible TPU skin are per-device. The collar shells are PETG for rigidity; the station back shell is
any rigid filament. The collar is assembled by seating the board and cell in the back shell, fitting
the PTFE membrane over the mic hole, bonding the shell halves with water-resistant glue, and slipping
the sealed unit into the TPU skin on a quick-release strap. The station is seated, glued closed, and
plugged in by the bowl. The hardware designs are released under CERN-OHL-S v2 (strongly reciprocal
open hardware).

\begin{figure}[ht]
  \centering
  \begin{subfigure}{0.40\textwidth}
    \centering\includegraphics[width=\linewidth]{collar-strap-front.jpg}
    \caption{Front, on the strap.}
    \label{fig:collar-front}
  \end{subfigure}\hfill
  \begin{subfigure}{0.40\textwidth}
    \centering\includegraphics[width=\linewidth]{collar-strap-back.jpg}
    \caption{Back, on the strap.}
    \label{fig:collar-back}
  \end{subfigure}
  \caption{The finished collar on a quick-release strap.}
  \label{fig:finished-collar}
\end{figure}

\subsection{Rationale for the split}

The collar must be small, light, and run for a useful time on a small cell, which rules out Wi-Fi on
the wearable: its power draw is incompatible with a collar-sized battery. BLE is low-power but
short-range and is not an internet transport. The fixed, mains-powered station resolves both
problems, giving the collar a cheap always-on BLE peer and a Wi-Fi path to the cloud. The split also
keeps the privacy boundary clean: the heavy and sensitive processing (including all audio) stays on
the collar, and only a classification result crosses to the station.
